\documentclass[11pt]{article}
\usepackage{fontspec}
\setmainfont{Times New Roman}
\setsansfont{Arial}
\setmonofont{Consolas}
\usepackage{amsmath,amssymb}
\usepackage{geometry}
\usepackage{microtype}
\geometry{a4paper,margin=3cm}
\setlength{\parindent}{1.25em}
\setlength{\parskip}{0pt}
\emergencystretch=18em
\allowdisplaybreaks
\raggedbottom
\providecommand{\vdotswithin}[1]{\mathrel{\vdots}}

\begin{document}

\section*{\S\ 5. Најобчејше области из целых трансцендентных чисел при основаньу на алгебраичној бази.}

Нехај $H$ буде која-коли алгебраична база всих чисел, и нехај
$\mathfrak M(\mathfrak G)$ означа совкупност всих областиј $\mathfrak G$, кторе
содержет $H$ и зато такоже $[H]$. Когда $H$ пробегаје всаку можну базу,
$\mathfrak M(\mathfrak G)$ пробегаје совкупност всих областиј $\mathfrak G$;
зато, како в \S\ 4, можемо ограничити се на разсмотрјенье областиј
$\mathfrak G$ из $\mathfrak M(\mathfrak G)$.

Области $\mathfrak G$ из $\mathfrak M(\mathfrak G)$ все содержет $[H]$ како
делитель; але, понеже само $[H]$ не јест област $\mathfrak G$, не можно, како в
\S\ 4, закльучити, же $[H]$ јест највечши обчи делитель. Же то всеже јест тако,
показује изчетна множина областиј $\mathfrak G$ из $\mathfrak M(\mathfrak G)$,
кторе взникајут малым обчньеньем модулној области из \S\ 3 и фактично не имајут
ниједного обчего елемента кроме $[H]$.

{\itshape Нехај области $\mathfrak G_\nu$ $(\nu=1,2,\ldots)$ за всако фиксовано
$\nu$ состоје из всих величин}
\begin{equation}
z=f(\eta)+g_1(\eta)\eta_1^\nu+g_2(\eta)\eta_2^\nu+\cdots+g_t(\eta)\eta_t^\nu,
\tag{1}
\end{equation}
{\itshape где $f(\eta)$ пробегаје все целочислене полиномы од $\eta$;
$\eta_1,\ldots,\eta_t$ пробегајут все базове величины, а
$g_1(\eta),\ldots,g_t(\eta)$ при всакој фиксованој комбинацији
$\eta_1^\nu,\ldots,\eta_t^\nu$ поједино пробегајут все целе алгебраичне функције
од $\eta$.}

Показује се аналогично како в \S\ 3, же областјам $\mathfrak G_\nu$ належет
абстрактне својства I--IV; $\mathfrak G_\nu$ содержи само целе алгебраичне
функције и при $\nu=1$ јест идентична с областју из \S\ 3. Ныне показујемо, же
области $\mathfrak G_\nu$ не имајут ниједного дальнего обчего елемента кроме
полиномов $f(\eta)$ из $[H]$. Именно, нехај $z$ буде која-коли цела
алгебраична, але не рационална, функција од $\eta$, ктора удовлетворјаје
уравненьу неразложимому взходно к $[H]$:
\begin{equation}
z^\chi+a_1(\eta)z^{\chi-1}+a_2(\eta)z^{\chi-2}+\cdots+a_\chi(\eta)=0,
\tag{2}
\end{equation}
где $\chi>1$; и нехај $\lambda$ буде највеча мед степенами полиномов
$a_1(\eta),\ldots,a_\chi(\eta)$. Ако $z$ належи к $\mathfrak G_\nu$, тогда
величина $\xi=z-f(\eta)$ удовлетворјаје такоже неразложимому, взходно к $[H]$,
уравненьу
\begin{equation}
\xi^\chi+b_1(\eta)\xi^{\chi-1}+b_2(\eta)\xi^{\chi-2}+\cdots+b_\chi(\eta)=0,
\tag{3}
\end{equation}
при чем $b_1(\eta),b_2(\eta),\ldots,b_\chi(\eta)$, како симетричне функције од
$\xi$, по (1) содержет мед членами најнижшего степена члене најманье степенов
$\nu,2\nu,\ldots,\chi\nu$. Сравньенье (2) и (3) даје
\[
 b_1(\eta)=a_1(\eta)+\chi f(\eta).
\]
Зато, ако величина $\xi=z+\frac{a_1(\eta)}{\chi}$ удовлетворјаје уравненьу
\begin{equation}
\xi^\chi+c_2(\eta)\xi^{\chi-2}+\cdots+c_\chi(\eta)=0,
\tag{4}
\end{equation}
тогда
\begin{equation}
 \xi-\zeta=\frac{b_1(\eta)}{\chi};
 \qquad
 c_i(\eta)\equiv b_i(\eta)\pmod{\frac{b_1(\eta)}{\chi}},
\tag{5}
\end{equation}
где $b_1(\eta)$ починаје членами најманье степена $\nu$. Ныне $c_i(\eta)$,
како показује сравньенье (2) с (4), сут степена $\chi$ в коефициентах
$a_i(\eta)$ из (2), и зато највыше степена $\chi\lambda$ в $\eta$; с другој
страны, все $b_i(\eta)$ починајут членами најманье степена $\nu$. Ако зато
изберемо $\nu>\chi\lambda$, из (5) следује
\[
 c_i(\eta)=0;\qquad
 \xi^\chi=0\quad\text{или}\quad
 \left(z+\frac{a_1(\eta)}{\chi}\right)^\chi=0,
\]
т. ј. $z$ противно предположеньу ставаје рационална функција од $\eta$. Зато
следује:

{\itshape Ако $z$ јест цела алгебраична, не рационална функција од $\eta$, тогда
$z$ не може належати к никакој области $\mathfrak G_\nu$ за
$\nu>\chi\lambda$, или:}

{\itshape Пресек всих областиј $\mathfrak G$ из $\mathfrak M(\mathfrak G)$ дан
јест областју целости $[H]$, ктора сама не јест област $\mathfrak G$; зато
$\mathfrak M(\mathfrak G)$ не јест затворјена взходно к твореньу пресеков.}

Зато такоже конструкција најобчејших областиј $\mathfrak G$ помочу алгебраичној
базы јест тежје прегледна неж конструкција областиј $\mathfrak H$. Области $[H]$
належет својства I, III и IV; зато, ако рационално-цело адјунгујемо којиколи
систем $\mathfrak S$, тогда за настајучу област целости $[H,\mathfrak S]$
својства I и III оставајут, меджутым како IV твори особно условје за
$\mathfrak S$. Да бы настала област $\mathfrak G$, треба такоже наложити II на
$\mathfrak S$ како условје, или другими словами: за всако полье $\mathfrak K$,
конечно над $(H)$\footnote{$(H)$ означа полье, составече се из всих рационалных
функциј од $\eta$ с алгебраичными числовыми коефициентами.}, мора обстајати
полье $\mathfrak L$ над $\mathfrak K$, конечно над $\mathfrak K$, тако же
$[H,\mathfrak S]$ содержи примитивны елемент из $\mathfrak L$.\footnote{Срав. к
тому \S\ 7.} Обратно, всака област $\mathfrak G$ из
$\mathfrak M(\mathfrak G)$ може такоже быти представльена в форме
$[H,\mathfrak S]$; зато достајемо все области $\mathfrak G$ из
$\mathfrak M(\mathfrak G)$, когда $\mathfrak S$ пробегаје все тако ограничене
системы. Структура системов $\mathfrak S$ покаже се простеје в следујучем
параграфу помочу «рационалној базы»; тым једночасно буде достата полна аналогија
с теоремами из \S\ 4.

\end{document}

